\section{Books Made in a Workshop}

\subsection{Not one book called Thomism}

Thomas wrote in forms whose purposes differ.  Biblical commentaries and
lecture reports follow sacred pages; Aristotelian expositions ordinarily seek
the philosopher's intention; disputed questions preserve the architecture of
academic determination; quodlibets answer issues proposed in public sessions;
polemical works intervene in live conflicts; the two \emph{Summae} order large
fields for different pedagogical purposes; sermons exhort; liturgical texts
are meant to be sung and prayed.  A sentence moved from one form to another may
lose its question, audience, level of authority, or stage of development.

Ancient catalogs, manuscript evidence, cross-references, style, and doctrine
have gradually separated authentic works from doubtful and spurious
attributions.  The Leonine Commission's critical edition remains a continuing
project.  Corpus Thomisticum offers an invaluable electronic corpus and
authentication aids, but an online text does not dissolve manuscript variants
or give every reportatio the status of an autograph.  The printed \emph{Opera
omnia} familiar to later readers contains layers that Thomas did not himself
assemble into one final system.

The chronology matters.  The \emph{Sentences} commentary belongs to a required
early exercise; the \emph{Summa contra gentiles} developed through several
years; the \emph{Summa theologiae} began after much teaching and remained
unfinished; positions can be refined.  A real change should be demonstrated
from controlled texts, not inferred whenever two translated phrases differ.
Conversely, harmony should not be manufactured by reading a late distinction
back into every early formulation.

\subsection{Hand, voice, and collaborators}

Thomas composed with quill, notebooks, exemplars, dictation, and human help.
Autograph material displays compressed handwriting and revision.  Early
witnesses describe him dictating to several secretaries, an image sometimes
turned into a feat beyond ordinary capacity.  The underlying collaborative
practice is credible and necessary: Reginald of Piperno and other socii
organized travel, books, oral delivery, copying, and perhaps posthumous
compilation.  Their service does not erase authorship; it makes authorship
material.

Reportationes raise a special limit.  A student's record may be the only form
in which a lecture or sermon survives.  It can witness Thomas's teaching while
also containing abbreviation, misunderstanding, rearrangement, or later
revision.  The source audit therefore does not use ``Thomas says'' uniformly
for an autograph, an author-supervised work, a reported lecture, an early
catalog attribution, and a later devotional text.

The workshop extended backward through translation.  Thomas read Aristotle in
Latin, often through changing versions; received Avicenna, Averroes, Maimonides,
and patristic sources through particular textual channels; and used florilegia
whose attributions could fail.  It extended forward through editors who
supplied titles, divisions, punctuation, parallel citations, and the
\emph{Supplement}.  His achievement is not diminished by this dependence.
Theologically, dependence is precisely how created intelligence works: it
receives, judges, and communicates goods it did not originate.

\subsection{Objection as an act of justice}

Thomas's mature article commonly begins with objections, introduces a contrary
authority, determines the question, and answers the objections.  The parts do
not all carry his final position.  Quoting an objection as ``Aquinas taught''
is a basic source error; quoting the \emph{sed contra} as though he endorsed
every premise can be one as well.  The determination and replies must be read
together.

At its best, the form disciplines disagreement.  It makes the strongest
located difficulty visible, distinguishes terms before resolving them, and
returns to what remains true in an opposing argument.  It does not guarantee
that Thomas possessed the best source, represented a social group from its own
voice, or escaped the blind spots of his setting.  Method can expose an
objection without repairing an unjust inherited premise.  A source-first
reader honors the intellectual virtue of the form by testing its evidence as
carefully as Thomas tested a syllogism.
